\subsection{通过其函数层刻画流形}

5.4.1. 设 X 是一个拓扑空间，Y 是一个集合。假设对于 X 的每个开集 U，给定一个从 U 到 Y 的映射集 \(\mathcal{L}(\mathrm{U})\)。如果函数族 \(\mathcal{L} = \{\mathcal{L}(\mathrm{U})\}\) 满足下述条件，我们称其为取值于 Y 的函数层：

设 \((\mathbf{U}_{i})_{i\in\mathbf{I}}\) 是 X 的开子集族，其并为 U，且 f 是从 U 到 Y 的映射；f 属于 \(\mathcal{L}(\mathbf{U})\) 的充分必要条件是：对于 I 中所有 i，\(f|\mathbf{U}_{i}\) 属于 \(\mathcal{L}(\mathbf{U}_{i})\)。

5.4.2. 设 X 和 Y 是两个流形；对于每个开集 \(U \subset X\)，令 \(\mathcal{L}(U)\) 为从 U 到 Y 的态射集；则 L 是取值于 Y 的函数层。

当 Y = K 时，将如此定义的层记为 \(C_{X}^{r}\)。

5.4.3. 设 X 是一个拓扑空间，S 是 Banach 空间的集合。对每个 \(E \in S\)，设 \(F_{E}\) 是 X 上取值于 E 的函数层。假设由 \(F_{E}\)（其中 \(E \in S\)）组成的族满足下述条件：

对于每个 \(x \in X\)，存在 x 的一个开邻域 U，以及一个空间

\(E_{0} \in S\)，以及一个把 U 映到 \(\varphi\) 的某个开子集上的同胚 \(E_{0}\)，使得对于每个开集 \(V \subset U\) 和每个 \(E \in S\)，集合 \(\mathcal{F}_{E}(V)\) 由函数 \(g \circ \varphi\) 组成，其中 g 遍历空间 \(\mathcal{C}^{r}(\varphi(V); E)\)。

于是，在 X 上存在一个 C＾\(C^{r}\)r 类且类型为 S 的流形结构，并且这种结构是唯一的；它与 X 上给定的拓扑相容，且在此结构下，\(F_{E}\) 是 X 上取值于 E 的 C＾\(C^{r}\)r 类函数层。

5.4.4. 设 X 是一个拓扑空间，F 是一个取值于 K 的函数层，满足下述条件：对于 X 的每个点 x，存在一个整数 n、x 的一个开邻域 U，以及一个把 U 映到  的某个开子集上的同胚 \(\varphi\)，使得对于 U 的每个开集 V，集合 \(K^{n}\)\(\mathcal{F}(V)\) 由函数 \(g \circ \varphi\) 组成，其中 g 遍历 \(C^{r}\) 的开集 \(\varphi(V)\) 上取值于 K 的 C＾\(K^{n}\)r 类函数集。

于是，在 X 上存在唯一的局部有限维 C＾\(C^{r}\)r 类流形结构，使得 \(F = C_{X}^{r}\)。

5.4.5. 设 X 和 X' 是两个局部有限维的 C＾\(C^{r}\)r 类流形，f 是从 X 到 X' 的连续映射。f 是态射的充分必要条件是：对于 X' 的每个开集 U' 以及每个函数 \(g \in \mathcal{C}^{r}(U'; K)\)，函数 \(g \circ f\) 属于 \(\mathcal{C}^{r}(U; K)\)，其中 \(U = f^{-1}(U')\)。

